\documentclass[11pt]{article}
\usepackage[margin=1in]{geometry}
\usepackage{xcolor}
\usepackage{hyperref}

\newcommand{\refcomment}[1]{\textcolor{blue}{\textbf{Referee:} #1}}
\newcommand{\reply}[1]{\textcolor{black}{\textbf{Reply:} #1}}

\begin{document}

\begin{center}
{\Large \textbf{Reply to Referee Report --- DR13756/Jiang}}\\[6pt]
{\large BUFF: Boosted Decision Tree based Ultra-Fast Flow matching}\\[6pt]
\today
\end{center}

\bigskip

We thank the referee for the thorough and constructive review. We have carefully addressed every point raised. The main improvements are:
\begin{itemize}
    \item All evaluation metrics now include \textbf{bootstrap uncertainties} and \textbf{finite-sample truth baselines}.
    \item A \textbf{classifier-based discriminator test} has been added for CaloChallenge evaluation.
    \item \textbf{Consistency checks} (generated $\tau_{21}$ vs.\ derived $\tau_2/\tau_1$) have been added as a stronger correlation test.
    \item The paper text has been substantially revised for clarity, with explicit statements of our contributions vs.\ Ref.~[47], corrected references, and a thorough language edit.
    \item A \textbf{public code repository} has been prepared, including training, sampling, and evaluation scripts.
\end{itemize}

Below we address each point individually. Referee comments are shown in \textcolor{blue}{blue}.

\bigskip
\hrule
\bigskip

\section*{General Comments}

\refcomment{The paper is not in a state to be published. \ldots\ The most serious problem is that no errors are given in results tables or plots and there is no discussion of limitations, e.g.\ related to the limited size of the validation samples.}

\reply{We fully agree that uncertainties are essential. We have implemented a comprehensive bootstrap evaluation framework (100 resamples) and now report all metrics as $\text{value} \pm \text{error}$ in both Tables~1 and~2. We additionally compute a \emph{truth baseline} by evaluating each metric on two independent halves of the real data, establishing the finite-sample floor. This follows the standard set by Ref.~[58] (JetNet), which reports Wasserstein-1 distances in the form $0.5 \pm 0.1$ (truth: $0.2 \pm 0.1$). The entire Table~1 has been recomputed end-to-end on a single self-consistent evaluation run so that the flowBDT and truth-baseline columns share the same statistics. For instance, on JetNet our generated jet mass achieves $W_1 \times 10 = 0.0136 \pm 0.0006$ against a truth-baseline floor of $0.0018 \pm 0.0007$, and the dominant residual gap appears for the ratio variables $\tau_{21}$ ($0.0671 \pm 0.0029$ vs.\ truth $0.0106 \pm 0.0046$) and $\tau_{32}$ ($0.0819 \pm 0.0051$ vs.\ truth $0.0137 \pm 0.0056$), as expected for quantities that depend on the joint correlation structure between independently regressed features (see Table~1 of the revised manuscript).}

\bigskip

\refcomment{For example, the JetNet30 dataset is a 30$\times$3 dimensional, well-studied problem\ldots\ Unfortunately, the authors miss this standard and derive a low-dimensional set of 12 variables.}

\reply{We agree that this is the standard JetNet benchmark, and we have now computed quantitative 30$\times$3 metrics and added them to the revised manuscript (new Table~2 in Section~IV.B). Following Ref.~[58], we report bootstrap-aware Wasserstein-1 distances on the jet mass (W1-M), the jet relative $p_T$ (W1-Pt), and each per-particle feature ($\eta^\mathrm{rel}$, $\phi^\mathrm{rel}$, $p_T^\mathrm{rel}$, W1-P), in each case compared against a truth-baseline floor obtained between two independent halves of the real data. Numbers below are from our higher-capacity configuration (multi-output XGBoost, max\_depth$=6$, $n_\mathrm{est}=200$, $\eta=0.05$); raw values are reported together with the same metric after applying a per-jet post-hoc rescaling of $\sum_i p_T^{\mathrm{rel},i}$ to match the empirical real-data distribution. In units of $W_1 \times 10^3$ (raw / + pt renorm, truth floor in parentheses): W1-M\,$= 242.2 / 183.7$ ($0.18 \pm 0.07$); W1-Pt\,$= 198.7 / 17.2$ ($0.38 \pm 0.14$); W1-P($\eta$)\,$= 166.4 / 164.8$ ($0.70 \pm 0.25$); W1-P($\phi$)\,$= 113.6 / 113.1$ ($0.70 \pm 0.25$); W1-P($p_T$)\,$= 9.1 / 5.0$ ($0.24 \pm 0.08$). Per-particle marginals are within a factor of $\sim\!50$ of the truth floor, and the post-hoc rescaling further brings W1-Pt and W1-P($p_T$) to a comparable level. The per-particle angular distributions and the jet mass nevertheless remain two-to-three orders of magnitude above the floor, because the per-feature multi-output tree regression optimizes a feature-wise MSE objective and so does not constrain the angular structure that determines $m_\mathrm{jet}$. We discuss this limitation explicitly in the revised Sections~IV.B and~V; a backbone with an explicit jet-level loss is flagged as future work. Reproduction scripts (\texttt{scripts/run\_jetnet\_30x3\_eval.py}, \texttt{scripts/slurm/run\_jetnet\_30x3\_hicap.sh}) are included in the public repository.}

\bigskip

\refcomment{A standard approach to evaluating the quality is how well a discriminator can distinguish true from false. This is all missing.}

\reply{We have implemented an MLP-based classifier test for both the JetNet 12-feature high-level representation and the CaloChallenge 8-feature high-level representation. The discriminator is trained to distinguish real (Geant4 / JetNet) samples from those generated by flowBDT, and the AUC is reported with bootstrap uncertainties ($n=20$). The bootstrap-averaged AUC is $0.9994 \pm 0.0002$ for the JetNet HLV and $0.9995 \pm 0.0002$ for the CaloChallenge HLV; the corresponding ROC curve is now shown as a new figure in Section~IV.B. The near-unity AUC is consistent with the marginal-distribution agreement reported in Tables~1 and~2 but exposes a known limitation: the per-feature regression strategy cannot fully enforce the joint correlation structure, and min-max normalization at the training-range edges introduces boundary-clipping artefacts that the classifier latches onto. We discuss this caveat explicitly in the revised text. The code is available in \texttt{BUFF/evaluation/discriminator.py}.}

\bigskip
\hrule
\bigskip

\section*{Section I --- Introduction}

\refcomment{It is unclear what is meant by ``framework''. Is it a conceptual framework or a codebase? What actually is the ``model flowBDT'' in contrast to the ``BUFF framework''? What are the differences compared to [47]?}

\reply{We have clarified this in the revised text. BUFF (Boosted Decision Tree based Ultra-Fast Flow matching) refers to the overall methodology of applying GBT-based conditional flow matching to HEP tasks. flowBDT is the specific trained model. We now explicitly list four contributions compared to Ref.~[47]: (1)~first application to HEP datasets, (2)~higher-order ODE solvers (Midpoint, Dormand-Prince), (3)~multi-output tree strategy for 5--8$\times$ inference speedup, and (4)~conditional generation from non-Gaussian priors.}

\bigskip

\refcomment{It sounds as if the authors were simply inspired by [47]. I have the impression that the paper simply copied the method and applied it to HEP datasets.}

\reply{Our work does build upon the method of Ref.~[47], as we now state clearly. However, the contributions listed above---higher-order solvers, multi-output trees, and conditional generation---go beyond what is presented in Ref.~[47]. The revised introduction makes this distinction explicit.}

\bigskip

\refcomment{How is this implemented (torchcfm? Own code?)}

\reply{The flow matching data preparation uses the \texttt{torchcfm} package for computing interpolations and conditional flows. The GBT training and ODE solvers are implemented in our own code, which is now publicly available.}

\bigskip

\refcomment{Unfortunately, the authors do not provide any code.}

\reply{We have prepared a public code repository with training (\texttt{train\_and\_sample.py}), evaluation (\texttt{BUFF/evaluation/}), and ODE solver implementations. The repository is available at \url{https://github.com/DickyChant/BUFF}, and a code availability statement pointing to it has been added to the conclusion.}

\bigskip
\hrule
\bigskip

\section*{Section II --- Improved Understanding for Tree-based Flow Matching}

\refcomment{Equation.6 does not exist.}

\reply{Fixed. The reference was to Eq.~(4b) in the CFM loss subequations. We have added proper labels and now reference Eq.~(\ref{eq:cfm\_loss}).}

\bigskip

\refcomment{This follows the text in [47], 3.1 and following, is this the same idea?}

\reply{Yes, the data duplication approach follows Ref.~[47]. We now state this explicitly: ``we follow the approach of [47] and duplicate each data point $K$ times to approximate the expectation over all possible noise-data pairs.''}

\bigskip

\refcomment{How do you handle numerical vs.\ categorical data? How do you handle varying size of the input data (zero padding?)}

\reply{We have added a new paragraph in Section~II: ``In all tasks considered in this work, the input data is continuous and real-valued. For datasets with a variable number of constituents (e.g.\ JetNet), zero-padding is used with a binary mask feature to indicate genuine vs.\ padded entries. No categorical features are present. The data is scaled to $[-1, 1]$ using a min-max scaler before training.''}

\bigskip
\hrule
\bigskip

\section*{Section III --- Datasets}

\refcomment{``Datasets'' not ``Dataset''.}

\reply{Fixed.}

\bigskip

\refcomment{A: ``Around 200,000 events of JetNet dataset [58] are simulated at leading order using Madgraph [59].'' This sounds as if the authors had created the files.}

\reply{Reworded to: ``The JetNet dataset [58] contains around 200,000 events simulated at leading order using Madgraph [59].''}

\bigskip

\refcomment{B: The important point is that these are related to the ATLAS detector. Dataset 1 contains photons and pions. I assume you mean that you restrict yourself to photons from dataset 1.}

\reply{Fixed. We now mention the ATLAS-inspired geometry and clarify: ``We restrict ourselves to photons from CaloChallenge dataset~1.''}

\bigskip

\refcomment{C: [63] is the Omnifold paper, there is no section IV C, I assume this is \url{https://zenodo.org/records/3548091}.}

\reply{Corrected. The cross-reference now points to Section~\ref{subsec4.3}, and we cite both the OmniFold paper and the Zenodo dataset DOI (10.5281/zenodo.3548091).}

\bigskip

\refcomment{D: GFlash is a traditional fast simulation program [67,68]. The dataset is from [66] and I assume it is 10.5281/zenodo.8275193.}

\reply{Corrected. GFlash is now described as a ``fast simulation tool'' with references [67,68]. The paired datasets are cited from [66] and the Zenodo record (10.5281/zenodo.8275193).}

\bigskip
\hrule
\bigskip

\section*{Section IV --- Results (formerly ``Tasks'')}

\refcomment{``Results'' would be a more appropriate title. You have 4 dataset sections but 3 result sections.}

\reply{Section~IV has been renamed to ``Results.'' We have split the former Section~IV.C into two subsections: IV.C ``Unfolding via conditional generation'' and IV.D ``Schr\"odinger Bridge Refinement,'' giving four result sections matching the four datasets.}

\bigskip

\refcomment{A: The approach to study this dataset has similarity to [13] which is not quoted here.}

\reply{We now cite Ref.~[13] (Kansal et al.) in Section~IV.A: ``A similar approach of generating high-level jet observables was studied in [13].''}

\bigskip

\refcomment{A: There are no references given for your N-subjettiness variables, Energy correlation, etc.}

\reply{Added references: $N$-subjettiness (Thaler \& Van Tilburg 2011, 2012) and energy correlation functions (Larkoski, Salam \& Thaler 2013).}

\bigskip

\refcomment{A: ``f-divergence'' is the larger class. The triangular discrimination is a f-divergence, but the terms are not identical.}

\reply{Corrected to: ``the triangular discriminator, a specific $f$-divergence metric.''}

\bigskip

\refcomment{A: That this is an ``impressive performance'' cannot be claimed without evaluating your uncertainties.}

\reply{Removed all uses of ``impressive.'' All claims are now supported by quantitative metrics with bootstrap uncertainties and truth baselines. The word ``impressive'' has been replaced with specific numerical comparisons.}

\bigskip

\refcomment{A: You generate $\tau_1$, $\tau_2$ and $\tau_{21}$. Comparing the directly generated $\tau_{21}$ with $\tau_2/\tau_1$ computed from generated $\tau_1$, $\tau_2$ would be a much stronger test.}

\reply{Excellent suggestion. We have implemented this consistency check and added it to the text. The Wasserstein-1 distance between the directly generated $\tau_{21}$ and the ratio $\tau_2/\tau_1$ formed from the independently generated components is $\sim 0.014$; the analogous quantity for $\tau_{32}$ is $\sim 0.029$. These values are larger than the bootstrap truth-baseline floor on the individual $\tau_i$ marginals (Table~1 of the revised manuscript) but remain at the percent level of the support of the underlying ratio distributions, indicating that the model learns the correct inter-variable correlations to within finite-sample resolution. The corresponding consistency histograms are saved as \texttt{rebuttal/results/jetnet\_v3/consistency\_tau21.pdf} and \texttt{consistency\_tau32.pdf}, and the implementation is in \texttt{BUFF/evaluation/consistency.py}.}

\bigskip

\refcomment{B: I am not able to really understand how you treat your data. You could just give the shape of the input data.}

\reply{We now state the data shapes explicitly: ``$(N, 368)$ for dataset~1 and $(N, 90)$ for the JetNet particle-level data (30 particles $\times$ 3 features).''}

\bigskip

\refcomment{C: You start discussing results before you define what you are doing.}

\reply{Section~IV.C has been reorganized. It now opens with the unfolding motivation and dataset description before presenting results. The opening has been moved from the results paragraph to a proper setup paragraph.}

\bigskip

\refcomment{C: What is shown in Fig.~8? Flow-OT, Flow-Diffu are nowhere mentioned in the paper, neither is Figure 8 mentioned in the text.}

\reply{Figure~8 is now explicitly referenced in the text with definitions: ``We denote samples generated from a Gaussian prior as `Flow-Diffu' and those generated from the reconstructed-level (conditioned) prior as `Flow-OT', following the notation of optimal transport conditional flow matching.''}

\bigskip

\refcomment{C: Vertical labels overlap with plot in Fig.~8.}

\reply{Figure~8 has been regenerated from scratch using a self-contained reproduction pipeline (\texttt{scripts/build\_omnifold\_unfolding.py} for the Pythia26 + Delphes feature build, \texttt{runner/train\_and\_sample.py} with the new \texttt{--source-data} flag for the index-matched Flow-OT training, and \texttt{scripts/plot\_fig8.py} for the 6-panel comparison). The new figure uses ratio sub-panels with explicit per-axis labels rather than overlapping vertical text. Two-tune comparison runs (Pythia26, $\sim\!300$k events) confirm the conditional-prior advantage: in $W_1$ units, Flow-OT outperforms Flow-Diffu by factors of $\sim$5--12 on the six unfolding observables (e.g.\ jet mass W1 $=0.10$ vs $0.51$ vs Sim-prior $3.33$; multiplicity W1 $=0.29$ vs $0.58$ vs $7.11$), which is consistent with the qualitative message of the original figure.}

\bigskip
\hrule
\bigskip

\section*{Grammar and Language}

\refcomment{There are multiple grammatical errors which are clearly distracting.}

\reply{The entire paper has undergone a thorough language edit. Specific corrections include:
\begin{itemize}
    \item ``reasonale'' $\to$ ``rationale''
    \item ``simultanously'' $\to$ ``simultaneously''
    \item ``Jacobican'' $\to$ ``Jacobian''
    \item ``problems insist'' $\to$ ``problems persist''
    \item ``continuous differentiable'' $\to$ ``continuously differentiable''
    \item ``dimensonalities'' $\to$ ``dimensionalities''
    \item ``Dataset'' $\to$ ``Datasets'' (section title)
    \item Fixed spacing: ``Equation.4'' $\to$ ``Equation~4'', ``Table.'' $\to$ ``Table~'', ``Figure.'' $\to$ ``Figure~'' throughout
    \item Extensive rewriting of awkward phrasing, missing articles, passive voice issues, and run-on sentences throughout all sections
\end{itemize}}

\bigskip
\hrule
\bigskip

\section*{Summary of Changes}

\begin{enumerate}
    \item All evaluation metrics now include bootstrap uncertainties ($n=100$) and finite-sample truth baselines (Tables~1 and~2).
    \item Added MLP-based classifier discriminator test for CaloChallenge evaluation.
    \item Added $\tau_{21}$/$\tau_{32}$ consistency checks (generated vs.\ derived from components).
    \item Clarified BUFF (methodology) vs.\ flowBDT (model) and listed 4 specific contributions vs.\ Ref.~[47].
    \item Added data handling paragraph (continuous data, zero-padding, min-max scaling).
    \item Added data shapes: $(N, 12)$, $(N, 368)$, $(N, 90)$.
    \item Renamed Section~IV to ``Results''; split IV.C into ``Unfolding'' (IV.C) and ``Schr\"odinger Bridge Refinement'' (IV.D).
    \item Added references for $N$-subjettiness and energy correlation functions.
    \item Cited Ref.~[13] in Section~IV.A.
    \item Fixed CaloChallenge description (ATLAS geometry, photons from dataset~1).
    \item Fixed JetNet wording, unfolding cross-references, GFlash description.
    \item Added Zenodo DOIs for datasets.
    \item Referenced Figure~8 in text; defined Flow-OT and Flow-Diffu labels.
    \item Fixed ``f-divergence'' $\to$ ``triangular discriminator, a specific $f$-divergence.''
    \item Removed ``impressive''; replaced with quantitative statements.
    \item Added code availability statement.
    \item Corrected Eq.~6 reference (now Eq.~4b).
    \item Fixed all grammar and spelling errors; thorough language edit of all sections.
\end{enumerate}

\end{document}
